\section{CoR-Fuse: Consensus-aware Routing Fusion}
\subsection{Problem Formulation}
For each slide \(n\), we denote the patch-level representation extracted by foundation model \(m\) as:
\begin{equation}
\mathbf{X}_m^{(n)}
=
[
\mathbf{x}_{m,1}^{(n)},
\ldots,
\mathbf{x}_{m,P_n}^{(n)}
]^\top
\in
\mathbb{R}^{P_n\times d_m},
\end{equation}
where \(P_n\) denotes the number of patches in slide \(n\), and \(d_m\) is the embedding dimension of encoder \(m\).
All encoders operate on coordinate-aligned patches extracted from the same WSI, enabling patch-level correspondence across heterogeneous representations.
Given a set of \(M\) foundation models \(\mathcal{E}=\{\mathbf{X}_1,\ldots,\mathbf{X}_M\}\), our goal is to learn a slide-level predictor:
\begin{equation}
p_\theta
(y^{(n)}
|
\mathbf{X}_1^{(n)},\ldots,\mathbf{X}_M^{(n)})
\end{equation}
The final prediction is obtained from the fused representation generated by the proposed fusion module.
\subsection{Two-Stage Optimization for Adaptive Routing}
To stabilize adaptive multi-encoder fusion, we adopt a two-stage optimization strategy consisting of mean-fusion warm start and consensus-guided adaptive routing.
CoR-Fuse performs routing at the slide level to capture semantic-level encoder utility rather than local texture patterns. 
The detailed discussion of slide-level routing design choices is provided in Appendix \ref{slide-level reasons}.
\subsubsection{Stage 1: Mean-Fusion Initialization}
Since pathology foundation models produce representations with different embedding dimensions, we first project all embeddings into a shared latent space:
\begin{equation}
    \mathbf{H}_{m}^{(n)}
    =
    P_m(\mathbf{X}_{m}^{(n)})
    \in
    \mathbb{R}^{P_n\times 512},
\end{equation}
where \(P_m(\cdot)\) denotes an encoder-specific projection head that projects the features into a shared 512-dimensional latent space.
We first train a reference fusion model using uniform mean fusion.
The fused patch representations are subsequently aggregated by an MIL classifier to produce slide-level predictions.
The reference model is trained by minimizing the cross-entropy loss:
\begin{equation}
    \mathcal{L}_{T}=\mathrm{CE}(p_\theta
(\mathbf{X}_1^{(n)},\ldots,\mathbf{X}_M^{(n)}), y^{(n)})
\end{equation}
where \(y^{(n)}\) denotes the ground-truth label of slide \(n\), and \(p_\theta(\cdot)\) denotes the predicted class probabilities produced by the reference model.
We use a strictly nested evaluation protocol to prevent information leakage from the outer test fold.
Within each outer training split, the checkpoint is selected based on inner-validation AUC and is subsequently frozen.
At the end of the warm-start stage, the projection heads and MIL classifier are transferred to initialize the adaptive fusion model.
\subsubsection{Stage 2: Consensus-Guided Adaptive Routing}
Given the initialized shared representation space, we introduce a lightweight encoder-level router to estimate slide-specific fusion weights.
Importantly, this pooling operation is only used for routing and does not replace the patch-level representations used by the downstream MIL classifier.
For each encoder \(m\), we aggregate its projected patch embeddings into a slide-level descriptor through mean pooling:
\begin{equation}
\mathbf{\mu}_{m}^{(n)}
=
\frac{1}{P_n}\sum_{p=1}^{P_n}\mathbf H_{m,p}^{(n)}\in\mathbb R^{512}.
\end{equation}
We then construct a consensus representation as:
\begin{equation}
\mathbf{\mu}_{con}^{(n)}
=
\frac{1}{M}
\sum_{m=1}^{M}
\mathbf{\mu}_{m}^{(n)}
\in
\mathbb{R}^{512}.
\end{equation}
The consensus representation summarizes the average semantic content across all encoders for that slide.
These consensus-conditioned encoder descriptors are subsequently used by the adaptive router to estimate slide-dependent routing scores.
Each encoder descriptor is concatenated with the consensus representation and fed into a two-layer MLP router:
\begin{equation}
    s_m^{(n)}=
    \mathrm{MLP}
    ([\mathbf{\mu}_{m}^{(n)},\mathbf{\mu}_{con}^{(n)}])
    +b_m,
\end{equation}
where \(b_m\) is a learnable encoder-specific bias capturing the global preference of each foundation model.
The routing scores \(s_m^{(n)}\) are converted into normalized weights \(w_m^{(n)}\) through softmax, and the final fused patch representation is obtained as:
\begin{equation}
    \mathbf{Z}_p^{(n)}=
    \sum_{m=1}^{M}w_m^{(n)}\mathbf{H}_{m,p}^{(n)}
\end{equation}
Although routing decisions are made at the slide level, fusion is performed on patch embeddings to preserve the spatial correspondence required by downstream MIL aggregation.
The router and fusion model are optimized jointly using only the downstream classification objective.
Unlike distillation-based routing approaches, our method does not require additional teacher-student matching objectives or attribution supervision.
\subsection{Fusion Prediction}
The routing weights are applied to the patch-level representations from all foundation models to construct the fused patch bag.
The fused representations are subsequently processed by an ABMIL classifier \citep{ilse2018abmil}, which performs patch-level attention aggregation and generates slide-level predictions.
The entire CoR-Fuse module is optimized with the same downstream classification objective.
